\section{Introduction}
\label{sec:intro}

\IEEEPARstart{S}{atellite} gravimetry from the Gravity Recovery and Climate Experiment (GRACE, 2002--2017) and its Follow-On (GRACE-FO, since 2018) measures month-to-month changes in Earth's gravity field. Over land, these changes are dominated by terrestrial water storage (TWS), the sum of groundwater, soil moisture, snow, surface water and glaciers. The record has shown regional freshwater trends worldwide~\cite{rodell2018emerging}, closed the sea-level budget together with altimetry and Argo~\cite{loomis2019regularization}, and traced ice-sheet mass loss across the change of mission~\cite{velicogna2020continuity}. It is now the only source of TWS as an Essential Climate Variable~\cite{rodell2023water}. Reviews of the mission's science and processing are given in~\cite{chen2022applications,rodell2023water}.

Three properties of the data limit how they can be used. The fields are coarse. The effective resolution is about 300\,km, and products delivered on 0.25$^\circ$ or 0.5$^\circ$ grids carry far less independent information than their pixel counts suggest~\cite{vishwakarma2018resolution,loomis2019regularization,bonin2026ocean}. The errors are structured. Spherical-harmonic solutions show north--south stripes from correlated errors, and the filters and regularisation that suppress the stripes also bias and leak the signal~\cite{kusche2007approximate,wiese2016leakage,loomis2019regularization}. And the record has holes: an eleven-month gap between the missions and scattered missing months before it~\cite{chen2022applications,landerer2020extending}.

Machine learning has been applied to all three. Convolutional networks learn the mismatch between land-surface models and GRACE~\cite{sun2019combining}. Bayesian networks bridge the inter-mission gap with uncertainty~\cite{mo2022bayesian}. Statistical and recurrent models reconstruct storage back to the early twentieth century~\cite{humphrey2019gracerec,yin2023gtwsmlrec,palazzoli2025graice}. Self-supervised encoder-decoders downscale GRACE to 0.5$^\circ$ globally~\cite{gou2024global}. Over the same period, diffusion models~\cite{ho2020ddpm,song2021scoresde} have become the dominant generative approach for downscaling, ensemble forecasting and data assimilation in atmospheric and ocean science~\cite{mardani2025corrdiff,price2025gencast,manshausen2025genda}. Their main appeal there is not lower point error but samples from a distribution: ensembles of spatially coherent fields whose spread can be calibrated.

This survey connects these literatures. Its questions are:
\begin{enumerate}
  \item What limits the information content of GRACE products, and what error models exist that a probabilistic method should respect?
  \item Which deep-learning methods have been applied to GRACE, at which stage of the processing chain, and how are they validated?
  \item Which diffusion-model techniques, developed in machine learning and in other geosciences, address the constraints that GRACE imposes: few training samples, no fine-scale truth, a non-trivial measurement operator, and a conservation requirement?
\end{enumerate}

\subsection{Review method}
\label{sec:method}

The literature was collected on 13 November 2026 by five parallel searches, one per sub-topic: GRACE measurement and processing; deep learning on GRACE products; machine learning in the gravimetry processing chain and complementary geodetic sensors; diffusion-model foundations; and diffusion and generative models in the Earth sciences. Sources were general web search, OpenAlex, Crossref, the arXiv API, the GFZ list of GRACE publications, Copernicus preprint servers and EGU abstracts. Each bibliographic record was checked against Crossref, OpenAlex or arXiv.

We apply one rule throughout: \emph{a work is cited only if its full text was obtained and read}. Text was extracted from each PDF, and every claim attributed to a paper here was checked against that text. Works known only from an abstract or a bibliographic record are not cited, however relevant they appear. This excludes a number of closed-access papers, and also open-access papers whose publishers blocked automated retrieval. Those are listed separately for retrieval and will be added in a later revision. The rule has consequences that the reader should keep in mind. Several canonical references for GRACE processing are absent here and are cited indirectly through the reviews that discuss them. The one diffusion-model study on GRACE products that our searches identified could only be read as an abstract and is therefore not discussed. Statements such as ``none of the methods reviewed'' refer to the works we read.

\subsection{Organisation}

Section~\ref{sec:grace} summarises the GRACE measurement system, its products, its effective resolution and its error models. Section~\ref{sec:mlgrace} reviews deep learning on gridded GRACE products by task. Section~\ref{sec:geodetic} covers machine learning upstream in the processing chain, complementary geodetic sensors and data assimilation. Section~\ref{sec:diffusion} introduces diffusion models with emphasis on conditioning and inverse problems. Section~\ref{sec:geodiff} reviews diffusion and other generative models in the Earth sciences. Section~\ref{sec:synthesis} maps GRACE tasks to transferable techniques and lists open problems. Section~\ref{sec:conclusion} concludes.
